\section{Introduction}
\label{sec:intro}

VegasCore's contribution is a machine-checked path from a protocol
program to the strategic games, programming-language guarantees, and
backend proof obligations it denotes.
This section motivates event graphs against ordinary game trees, works
through a small Matching-Pennies example, sketches the semantic pipeline,
explains where runtime backends fit relative to the graph machine, and
then summarizes the main theorem surface.  The notation and glossary
immediately after the introduction define recurring terms such as
frontier, checkpoint, and kernel game.

\Cref{fig:overview} is a roadmap of the whole development: a protocol
program is compiled to a visibility-aware event graph, which executes on a
schedule-agnostic operational machine and presents as the strategic games on
which equilibrium and information properties are proved; concrete runtimes
enter from below as a composable tower of refinements that preserve the game's
utilities, bottoming out at the on-chain execution that is the end goal.

\begin{figure}[t]
\centering
\begin{tikzpicture}[
  >={Stealth[length=2.4mm]},
  font=\small,
  box/.style={draw, rounded corners, align=center, inner sep=4pt,
              minimum height=9mm, text width=46mm},
  core/.style={box, fill=blue!8, very thick},
  side/.style={box, fill=green!8, text width=42mm},
  goal/.style={box, dashed, fill=black!4, text width=46mm},
  flow/.style={->, thick},
  lab/.style={font=\footnotesize, midway, right=1mm, align=left},
]
% ── forward semantic spine (top-down) ────────────────────────────────
\node[box]  (prog)  at (0,0)    {Protocol program\\{\footnotesize public samples \(\cdot\) sealed commits \(\cdot\) reveals \(\cdot\) payoffs}};
\node[core] (graph) at (0,-2.1) {Visibility-aware \textbf{event graph}\\{\footnotesize the semantic center}};
\node[box]  (mach)  at (0,-4.2) {Operational machine\\{\footnotesize schedule-agnostic execution}};
\draw[flow] (prog)  -- node[lab]{compile\\{\scriptsize source syntax erased}} (graph);
\draw[flow] (graph) -- node[lab]{execute} (mach);
% ── strategic analysis surface (to the right) ────────────────────────
\node[side] (games) at (7.2,-1.6) {Strategic games\\{\footnotesize pure \(\cdot\) behavioral \(\cdot\) mixed-pure}};
\node[side] (pres)  at (7.2,-3.6) {FOSG / EFG presentations};
\draw[flow] (mach.east) .. controls +(2.2,0) and +(-1.6,0) .. (games.west)
  node[pos=0.55,above,font=\footnotesize,align=center]{frontier\\view};
\draw[flow] (games) -- (pres);
\node[align=left, text width=46mm, font=\footnotesize] at (7.2,0.4)
  {\textbf{Prove \& analyze:} payoff adequacy, visibility, Kuhn equivalence,
   Nash / equilibria};
% ── runtime tower (bottom-up refinement) ─────────────────────────────
\node[box]  (back) at (0,-6.4) {Abstract backend layers\\{\footnotesize audited steps \(\cdot\) encoded storage \(\cdot\) message buffers}};
\node[goal] (evm)  at (0,-8.6) {\textbf{EVM / smart-contract runtime}\\{\footnotesize end goal: on-chain execution}};
\draw[flow] (back) -- node[lab]{refines\\{\scriptsize utility-preserving}} (mach);
\draw[flow] (evm)  -- node[lab]{refines\\{\scriptsize layers compose into a tower}} (back);
% ── legend ───────────────────────────────────────────────────────────
\node[draw, rounded corners, align=left, font=\scriptsize, fill=white,
      inner sep=3pt, anchor=north west] at (4.6,-6.0)
  {\tikz\node[draw,rounded corners,minimum width=4mm,minimum height=3mm,fill=blue!8]{};\ mechanized in this work\\
   \tikz\node[draw,rounded corners,dashed,minimum width=4mm,minimum height=3mm,fill=black!4]{};\ motivating target};
\end{tikzpicture}
\caption{VegasCore at a glance.  Solid boxes are mechanized in this
development.  The downward arrows are the forward semantics that turn a
protocol into a game; the strategic surface on the right is where
game-theoretic properties are proved.  Concrete runtimes enter from below as a
\emph{tower} of utility-preserving refinements that compose down to the
primitive machine; the EVM runtime (dashed) is the motivating endpoint the
interfaces are designed to reach.}
\label{fig:overview}
\end{figure}

\subsection{Protocols vs. game trees}

Many finite games are naturally written as protocols.  Each player
commits to a private value, a public sample is drawn, hidden values are
revealed, and the terminal public state determines the payoffs.  Writing
the same object directly as an
extensive-form tree forces an arbitrary linearization of events that are
intended to be simultaneous or independent.  The tree can repair that
choice with information sets, but the independence structure is then
spread across syntax and equivalence classes.

VegasCore takes the dependency structure as primary.  A program is
checked and compiled to a finite event graph.  The graph records which
fields each event reads, which field each event writes, and which earlier
events are prerequisites because they provide information needed at the
event.  The graph state records only completed nodes and the typed store.
It does not record the order in which independent nodes happened to
complete.

The result is a schedule-agnostic operational semantics.  Primitive
execution steps one available event at a time, but the scheduler is an
external event law.  The strategic semantics use a frontier round view
(a view of the currently ready graph boundary): players see checkpoint
observations and choose for ready commit nodes; internal graph work is
closed by the checkpoint policy.  This separates the schedule-agnostic
game events of the language from implementation-level linearization
detail.

That separation is why the formal development factors protocol meaning
through an event graph rather than generating an EFG directly.

\subsection{A running example}
\label{sec:intro:matching-pennies}

Consider a Matching-Pennies protocol.  The matcher and mismatcher each
commit a hidden boolean, then both booleans are revealed.  The matcher
gets \(+1\) on equality and \(-1\) otherwise; the mismatcher gets the
opposite payoff.  A source program necessarily lists the constructors in
some order, for example:
\begin{quote}\small
\begin{tabular}{@{}l@{}}
\(\Commit_{\mathsf{m}}\,x_{\mathsf{m}} \mid \mathsf{true};\)\\
\(\Commit_{\mathsf{m'}}\,x_{\mathsf{m'}} \mid \mathsf{true};\)\\
\(\Reveal\,y_{\mathsf{m}} := x_{\mathsf{m}};\)\\
\(\Reveal\,y_{\mathsf{m'}} := x_{\mathsf{m'}};\)\\
\(\Ret\,[\mathsf{m}\mapsto u(y_{\mathsf{m}},y_{\mathsf{m'}}),
          \mathsf{m'}\mapsto -u(y_{\mathsf{m}},y_{\mathsf{m'}})].\)
\end{tabular}
\end{quote}
The intended game is not sequential at the commitment phase.  The two
commits are independent choices because neither player can observe the
other player's hidden field before committing.

\begin{figure}[t]
\centering
\begin{tikzpicture}[
  every node/.style={font=\small},
  ev/.style={draw, rounded corners, inner sep=2pt, minimum width=1.8cm,
             minimum height=4mm, fill=blue!6},
  arr/.style={-{Stealth[length=2mm]}, thick},
]
\node[ev] (cm) at (-2,1) {\(\Commit_{\mathsf{m}}\)};
\node[ev] (cmp) at (2,1) {\(\Commit_{\mathsf{m'}}\)};
\node[ev] (rm) at (-2,-1) {\(\Reveal_{\mathsf{m}}\)};
\node[ev] (rmp) at (2,-1) {\(\Reveal_{\mathsf{m'}}\)};
\node[ev] (ret) at (0,-2.5) {\(\Ret\)};
\draw[arr] (cm) -- (rm);
\draw[arr] (cmp) -- (rmp);
\draw[arr] (cm) -- (rmp);
\draw[arr] (cmp) -- (rm);
\draw[arr] (rm) -- (ret);
\draw[arr] (rmp) -- (ret);
\end{tikzpicture}

\caption{Matching Pennies as an event graph.  Straight edges from a
commit to its own reveal and from reveals to \(\Ret\) are data
dependencies.  Cross-edges from one player's commit to the other
player's reveal are barrier dependencies: they enforce that all earlier
commitments are fixed before any reveal phase begins, but they carry no
hidden payload.}
\label{fig:matching-pennies-eg}
\end{figure}

In \Cref{fig:matching-pennies-eg}, the two commit nodes share the
initial frontier, the graph boundary of currently ready events.  Each
reveal depends on its own hidden source field, and the reveal phase also
waits for prior commits.  This second dependency is the commit/reveal
barrier: a reveal must not become available while an earlier commitment
is still mutable.  The barrier is graph structure, not a strategy-side
convention.  Thus the source text may list one commitment before the
other, but the compiled game still presents them as simultaneous hidden
choices.

\subsection{The semantic spine}
\label{sec:intro:spine}

The construction has the following shape:
\[
\begin{aligned}
  \mathsf{VegasLang}
  &\xrightarrow{\;\mathsf{lower}\;}
  \mathsf{VegasCore}
  \xrightarrow{\;\mathsf{compile}\;}
  \mathsf{EventGraph} \\
  &\xrightarrow{\;\mathsf{primitiveMachine}\;}
  \mathsf{Machine}
  \xrightarrow{\;\mathsf{frontierRoundView}\;}
  \mathsf{KernelGame}.
\end{aligned}
\]
\(\mathsf{VegasLang}\) is a surface layer over the concrete expression
language.  It contains administrative \(\Let\)-bindings and nullable
yield sugar.  Lowering erases administrative lets by substitution and
turns yields into a hidden optional commit followed by a reveal.  The
maintained core, \(\mathsf{VegasCore}\), contains only protocol
constructors: \(\Sample\), \(\Commit\), \(\Reveal\), and \(\Ret\).

Compilation from a \(\mathsf{GraphProgram}\) produces a
\(\mathsf{CompiledProgram}\): a graph, a proof that the graph is
well-formed, and terminal payoff projections.  Compilation does not
need the stronger checked-game obligations.  A \(\mathsf{WFProgram}\)
adds reveal completeness and legal commit guards; those obligations are
used later to prove guard liveness, checkpoint progress, and the
game-facing theorem surface.

The event graph is source-independent.  It is list-backed in the Lean
implementation, but source syntax is not graph identity: graph fields
are the initial fields followed by one output field per event row, and
node ids are canonical numeric positions in the node list.  Two
alpha-renamed surface programs can compile to equal graph data when
their protocol structure is equal.

The final arrow to \(\mathsf{KernelGame}\) is where strategic analysis
starts.  A kernel game does not remember the syntax of the protocol; it
remembers the strategy spaces, the probability law over outcomes, and
the utility function.  The source-payoff adequacy theorems prove that
the outcomes used there are exactly the source payoffs evaluated at
terminal graph states.

\subsection{Where runtime backends fit}
\label{sec:intro:runtime-backends}

Concrete runtimes enter below the primitive event-graph machine.  A
smart-contract runtime, a payment-channel state machine, or a local
simulator is modeled as a lower-level machine: its states contain
runtime storage and administrative data, its events model the relevant
calls or operational steps, and its observations say what each player can
infer from that runtime.  To connect such a machine to VegasCore, the
backend proves a stochastic refinement to the primitive event-graph
machine and supplies a trace law showing that the implementation realizes
the selected frontier game surface.  Ethereum/EVM is one motivating
instance of this backend pattern.

This is intentionally stronger than saying that a contract ``looks like''
the protocol on examples: the refinement proof identifies which
implementation details are erased, which events correspond to graph
events, and why the final utilities are unchanged.  Extra communication
requires a separate strategy-space argument; \Cref{sec:ref:cheap-talk}
gives the three sound positions.

The backend layer is a tower-building interface: stochastic refinements
compose, event-batch law families compose with them, and abstract backend
examples demonstrate audited administrative steps, encoded storage, and
explicit message buffers.  The message-in-flight wrapper is deliberately
player-visible, so it can model settings in which strategies condition
on pending messages.  When those messages are irrelevant,
inert-extension theorems transport equilibria; when they are
strategically meaningful, the specification game includes the
corresponding protocol messages or cheap-talk channel.

\subsection{Contributions and theorem surface}
\label{sec:intro:contributions}

The main contribution is the checked semantic path from protocol syntax
to game presentations and runtime-refinement obligations.  The public
theorem surface is organized around the following claims.
\begin{enumerate}[leftmargin=1.4em,itemsep=2pt]
\item \textbf{Compilation.}  Checked programs compile to well-formed,
  source-free event graphs.  Payoff projections read only public fields
  that are available by graph termination.
\item \textbf{Progress.}  Source guard legality compiles to graph
  \(\mathsf{GuardLive}\), and graph guard liveness gives progress for
  the primitive downset checkpoint model.
\item \textbf{Outcome adequacy.}  Terminal primitive-machine outcomes
  agree with the source payoff expressions reconstructed from the
  terminal compiled store.
\item \textbf{Visibility.}  Commit reads are exactly the committing
  player's visible dependency footprint; public samples and payoff
  projections read declared public dependencies; reveals are the only
  hidden-to-public flow.
\item \textbf{Frontier semantics.}  Pure, behavioral, and product
  mixed-pure frontier games are defined from the canonical frontier
  round view.  Their outcome kernels are source-payoff projections of
  the same canonical run distribution.
\item \textbf{FOSG and EFG presentations.}  The FOSG and EFG layers are
  payoff-faithful presentations of the frontier semantics.
\item \textbf{Kuhn.}  Behavioral and product mixed-pure frontier profiles
  are related by a strategy/deviation equivalence at the canonical
  round-view information model: outcome laws agree, and one-player
  deviations can be matched across the two presentations.
\item \textbf{Runtime refinement.}  A stochastic refinement plus an
  adequate strategy-indexed trace law preserves the utility distribution
  of the selected frontier game surface.  Refinements and law-family
  lifts compose, so a tower of abstract backend layers can be collapsed
  to one proof obligation against the primitive machine.  With the
  strategy-space conditions of \Cref{sec:ref:cheap-talk} and bounded
  utility for the compared trace utilities, Nash equilibria on that
  surface pull back to the implementation trace game.  The direction is
  one-way: implementation equilibria do not automatically project back to
  specification equilibria.
\end{enumerate}

\leancorr{The paper-facing theorem index is
\TT{Vegas.Core.Theorems.Claims}.
The construction modules are \TT{Vegas.Lang}, \TT{Vegas.Core},
\TT{Vegas.EventGraph}, \TT{Vegas.Core.ToEventGraph},
\TT{Vegas.GameBridge}, \TT{Vegas.Export}, and \TT{Vegas.Machine}.}
